% GENERATED FILE. Edit canonical lessons, then run npm run generate:books.
% canonical-source-hash: fnv1a64:bd4532d0f452a815
% canonical-lessons: PA-C188-meeting-return, PA-C188-milna, PA-C188-milna-review, PA-C188-milde, PA-C188-milna-milde-contrast, PA-C188-milde-review, PA-C188-milia, PA-C188-family-contrast, PA-C188-milia-review, PA-C188-payoff

\chapter{To Meet; We Meet; I Met}
\label{ch:pa-a2-meeting-family}
\hlchaptermodality{188}

\begin{chapteropening}
\textbf{By the end of this chapter:} \emph{I can distinguish and use a meeting infinitive, a regular-group meeting form, and one completed masculine-singular meeting frame.}

Retrieve the bounded infinitive, habitual-plural, and masculine-singular perfective meeting contrast through listening, speaking, reading, and writing.
\end{chapteropening}

\section[\texorpdfstring{\pa{ਫਿਰ} \pa{ਮਿਲਾਂਗੇ} (phir milaange)}{ਫਿਰ ਮਿਲਾਂਗੇ (phir milaange)}]{Start with the meeting farewell you know}

\label{lesson:PA-C188-meeting-return}



\begin{quote}
\emph{Hear:} \emph{phir milaange}; YOU SAY: \textbf{\pa{ਫਿਰ} \pa{ਮਿਲਾਂਗੇ}}
\end{quote}

\begin{grammarlens}[title={Grammar Lens}]
Keep the familiar idea of meeting steady. The chapter will separate a naming form, a regular-group form, and one completed masculine-singular form, one owner lesson at a time.
\end{grammarlens}

\subsection*{Your turn}
\begin{itemize}
\raggedright
  \item \emph{Hear:} \emph{phir milaange}; YOU SAY: \textbf{\pa{ਫਿਰ} \pa{ਮਿਲਾਂਗੇ}}
  \item \emph{Read it:} \textbf{\pa{ਫਿਰ} \pa{ਮਿਲਾਂਗੇ}}
\end{itemize}

\begin{tcolorbox}[breakable,colback=teal!4,colframe=teal!35!black,title={Before you move on}]
Say the familiar farewell once more. This voice-first return owns no new form.
\end{tcolorbox}
\section[\texorpdfstring{\pa{ਮਿਲਣਾ} (milnaa)}{ਮਿਲਣਾ (milnaa)}]{The infinitive for “to meet”}

\label{lesson:PA-C188-milna}



\begin{quote}
Retrieve \textbf{\pa{ਫਿਰ} \pa{ਮਿਲਾਂਗੇ}} without looking.
\end{quote}

\begin{grammarlens}[title={Grammar Lens: name the action}]
\begin{quote}
\textbf{\pa{ਮਿਲਣਾ}} — \emph{milnaa} — to meet
\end{quote}

Treat this as one whole infinitive. It names the action; later lessons will contrast it with a regular-group form and a completed form.
\end{grammarlens}

\subsection*{Your turn}
\begin{itemize}
\raggedright
  \item \emph{Hear:} \emph{milnaa}; YOU SAY: \textbf{\pa{ਮਿਲਣਾ}}
  \item \emph{Read it:} \textbf{\pa{ਮਿਲਣਾ}}
  \item \emph{Write it:} copy \textbf{\pa{ਮਿਲਣਾ}} once
\end{itemize}

\begin{tcolorbox}[breakable,colback=teal!4,colframe=teal!35!black,title={Before you move on}]
Give the infinitive meaning “to meet.”
\end{tcolorbox}
\section[\texorpdfstring{\pa{ਮਿਲਣਾ} (milnaa)}{ਮਿਲਣਾ (milnaa)}]{Retrieve the infinitive}

\label{lesson:PA-C188-milna-review}



\begin{quote}
Say “to meet” before uncovering the line.
\end{quote}

\begin{grammarlens}[title={Grammar Lens}]
\textbf{\pa{ਮਿਲਣਾ}} names the action. Keep that job distinct from the new form in the next lesson.
\end{grammarlens}

\subsection*{Your turn}
\begin{itemize}
\raggedright
  \item \emph{Hear:} \emph{milnaa}; YOU SAY: \textbf{\pa{ਮਿਲਣਾ}}
  \item \emph{Write it:} cover the model and write \textbf{\pa{ਮਿਲਣਾ}}
\end{itemize}

\begin{tcolorbox}[breakable,colback=teal!4,colframe=teal!35!black,title={Before you move on}]
Read \textbf{\pa{ਮਿਲਣਾ}} once and state its job: it names the action.
\end{tcolorbox}
\section[\texorpdfstring{\pa{ਅਸੀਂ} \pa{ਮਿਲਦੇ} \pa{ਹਾਂ} (asiin milde haan)}{ਅਸੀਂ ਮਿਲਦੇ ਹਾਂ (asiin milde haan)}]{A regular meeting by a group}

\label{lesson:PA-C188-milde}



\begin{quote}
Retrieve \textbf{\pa{ਮਿਲਣਾ}}, then the familiar words \textbf{\pa{ਅਸੀਂ}} and \textbf{\pa{ਹਾਂ}}.
\end{quote}

\begin{grammarlens}[title={Grammar Lens: the regular-group form}]
\begin{quote}
\textbf{\pa{ਅਸੀਂ} \pa{ਮਿਲਦੇ} \pa{ਹਾਂ}} — \emph{asiin milde haan} — we meet
\end{quote}

The middle word is the measured new form. In this bounded frame it presents a regular or habitual meeting by a plural group.
\end{grammarlens}

\subsection*{Your turn}
\begin{itemize}
\raggedright
  \item \emph{Hear:} \emph{asiin milde haan}; YOU SAY: \textbf{\pa{ਅਸੀਂ} \pa{ਮਿਲਦੇ} \pa{ਹਾਂ}}
  \item \emph{Read it:} \textbf{\pa{ਅਸੀਂ} \pa{ਮਿਲਦੇ} \pa{ਹਾਂ}}
  \item \emph{Write it:} copy \textbf{\pa{ਮਿਲਦੇ}} once
\end{itemize}

\begin{tcolorbox}[breakable,colback=teal!4,colframe=teal!35!black,title={Before you move on}]
Say the regular-group frame once without looking.
\end{tcolorbox}
\section[\texorpdfstring{\pa{ਮਿਲਣਾ} · \pa{ਅਸੀਂ} \pa{ਮਿਲਦੇ} … (milnaa; asiin …)}{ਮਿਲਣਾ · ਅਸੀਂ ਮਿਲਦੇ … (milnaa; asiin …)}]{Name the action or describe a regular meeting}

\label{lesson:PA-C188-milna-milde-contrast}



\begin{quote}
Say the naming form, then the regular-group form.
\end{quote}

\begin{grammarlens}[title={Grammar Lens: two jobs}]
\begin{itemize}
\raggedright
  \item \textbf{\pa{ਮਿਲਣਾ}} — names “to meet.”
  \item \textbf{\pa{ਅਸੀਂ} \pa{ਮਿਲਦੇ} \pa{ਹਾਂ}} — says “we meet” as a regular pattern.
\end{itemize}
\end{grammarlens}

\subsection*{Your turn}
\begin{itemize}
\raggedright
  \item \emph{Hear:} \emph{milnaa}; YOU SAY: \textbf{\pa{ਮਿਲਣਾ}}
  \item \emph{Hear:} \emph{asiin milde haan}; YOU SAY: \textbf{\pa{ਅਸੀਂ} \pa{ਮਿਲਦੇ} \pa{ਹਾਂ}}
  \item \emph{Write it:} \textbf{\pa{ਮਿਲਣਾ}} and \textbf{\pa{ਮਿਲਦੇ}} under the correct English jobs
\end{itemize}

\begin{tcolorbox}[breakable,colback=teal!4,colframe=teal!35!black,title={Before you move on}]
Produce both forms in meaning order: action name, regular group.
\end{tcolorbox}
\section[\texorpdfstring{\pa{ਅਸੀਂ} \pa{ਮਿਲਦੇ} \pa{ਹਾਂ} (asiin milde haan)}{ਅਸੀਂ ਮਿਲਦੇ ਹਾਂ (asiin milde haan)}]{Retrieve the regular-group frame}

\label{lesson:PA-C188-milde-review}



\begin{quote}
Hear “we meet regularly” and prepare the complete frame.
\end{quote}

\begin{grammarlens}[title={Grammar Lens}]
Do not replace the regular-group form with the infinitive. Read the subject, the meeting form, and the copula as one measured frame.
\end{grammarlens}

\subsection*{Your turn}
\begin{itemize}
\raggedright
  \item \emph{Hear:} \emph{asiin milde haan}; YOU SAY: \textbf{\pa{ਅਸੀਂ} \pa{ਮਿਲਦੇ} \pa{ਹਾਂ}}
  \item \emph{Read it:} \textbf{\pa{ਅਸੀਂ} \pa{ਮਿਲਦੇ} \pa{ਹਾਂ}}
  \item \emph{Write it:} cover and write \textbf{\pa{ਅਸੀਂ} \pa{ਮਿਲਦੇ} \pa{ਹਾਂ}}
\end{itemize}

\begin{tcolorbox}[breakable,colback=teal!4,colframe=teal!35!black,title={Before you move on}]
Give the middle form alone, then restore the full frame.
\end{tcolorbox}
\section[\texorpdfstring{\pa{ਮੈਂ} \pa{ਅਮਨ} \pa{ਨੂੰ} \pa{ਮਿਲਿਆ} (main Aman nuun miliaa)}{ਮੈਂ ਅਮਨ ਨੂੰ ਮਿਲਿਆ (main Aman nuun miliaa)}]{One completed meeting}

\label{lesson:PA-C188-milia}



\begin{quote}
Retrieve \textbf{\pa{ਮੈਂ}}, \textbf{\pa{ਅਮਨ}}, and \textbf{\pa{ਨੂੰ}}, then the action name and the regular-group meeting form from the last lessons.
\end{quote}

\begin{grammarlens}[title={Grammar Lens: masculine singular completed form}]
\begin{quote}
\textbf{\pa{ਮੈਂ} \pa{ਅਮਨ} \pa{ਨੂੰ} \pa{ਮਿਲਿਆ}} — \emph{main Aman nuun miliaa} — I met Aman
\end{quote}

The final word is the measured masculine-singular completed form. This chapter does not generalize beyond that bounded agreement frame.
\end{grammarlens}

\subsection*{Your turn}
\begin{itemize}
\raggedright
  \item \emph{Hear:} \emph{main Aman nuun miliaa}; YOU SAY: \textbf{\pa{ਮੈਂ} \pa{ਅਮਨ} \pa{ਨੂੰ} \pa{ਮਿਲਿਆ}}
  \item \emph{Read it:} \textbf{\pa{ਮੈਂ} \pa{ਅਮਨ} \pa{ਨੂੰ} \pa{ਮਿਲਿਆ}}
  \item \emph{Write it:} copy \textbf{\pa{ਮਿਲਿਆ}} once
\end{itemize}

\begin{tcolorbox}[breakable,colback=teal!4,colframe=teal!35!black,title={Before you move on}]
Say the completed frame once without looking.
\end{tcolorbox}
\section[\texorpdfstring{\pa{ਮਿਲਣਾ} · \pa{ਅਸੀਂ} \pa{ਮਿਲਦੇ} … (milnaa; asiin …)}{ਮਿਲਣਾ · ਅਸੀਂ ਮਿਲਦੇ … (milnaa; asiin …)}]{Contrast the three meeting jobs}

\label{lesson:PA-C188-family-contrast}



\begin{quote}
Prepare three meanings: to meet, we meet regularly, I met Aman.
\end{quote}

\begin{grammarlens}[title={Grammar Lens: one family, three bounded jobs}]
\begin{itemize}
\raggedright
  \item \textbf{\pa{ਮਿਲਣਾ}} — names the action.
  \item \textbf{\pa{ਅਸੀਂ} \pa{ਮਿਲਦੇ} \pa{ਹਾਂ}} — gives the regular-group frame.
  \item \textbf{\pa{ਮੈਂ} \pa{ਅਮਨ} \pa{ਨੂੰ} \pa{ਮਿਲਿਆ}} — gives the completed masculine-singular frame.
\end{itemize}
\end{grammarlens}

\subsection*{Your turn}
\begin{itemize}
\raggedright
  \item \emph{Say it:} \textbf{\pa{ਮਿਲਣਾ}} for “to meet”
  \item \emph{Say it:} \textbf{\pa{ਅਸੀਂ} \pa{ਮਿਲਦੇ} \pa{ਹਾਂ}} for the regular pattern
  \item \emph{Say it:} \textbf{\pa{ਮੈਂ} \pa{ਅਮਨ} \pa{ਨੂੰ} \pa{ਮਿਲਿਆ}} for the completed meeting
  \item \emph{Write it:} \textbf{\pa{ਮਿਲਣਾ} · \pa{ਮਿਲਦੇ} · \pa{ਮਿਲਿਆ}} in meaning order
\end{itemize}

\begin{tcolorbox}[breakable,colback=teal!4,colframe=teal!35!black,title={Before you move on}]
Cover the models and produce the same three-way contrast.
\end{tcolorbox}
\section[\texorpdfstring{\pa{ਮੈਂ} \pa{ਅਮਨ} \pa{ਨੂੰ} \pa{ਮਿਲਿਆ} (main Aman nuun miliaa)}{ਮੈਂ ਅਮਨ ਨੂੰ ਮਿਲਿਆ (main Aman nuun miliaa)}]{Retrieve the completed meeting}

\label{lesson:PA-C188-milia-review}



\begin{quote}
Hear “I met Aman” and prepare the bounded frame.
\end{quote}

\begin{grammarlens}[title={Grammar Lens}]
Keep the completed form distinct from the action name and the regular-group form. No broader agreement claim is being made here.
\end{grammarlens}

\subsection*{Your turn}
\begin{itemize}
\raggedright
  \item \emph{Hear:} \emph{main Aman nuun miliaa}; YOU SAY: \textbf{\pa{ਮੈਂ} \pa{ਅਮਨ} \pa{ਨੂੰ} \pa{ਮਿਲਿਆ}}
  \item \emph{Read it:} \textbf{\pa{ਮੈਂ} \pa{ਅਮਨ} \pa{ਨੂੰ} \pa{ਮਿਲਿਆ}}
  \item \emph{Write it:} cover and write \textbf{\pa{ਮੈਂ} \pa{ਅਮਨ} \pa{ਨੂੰ} \pa{ਮਿਲਿਆ}}
\end{itemize}

\begin{tcolorbox}[breakable,colback=teal!4,colframe=teal!35!black,title={Before you move on}]
Give the completed form alone, then restore the full frame.
\end{tcolorbox}
\section[\texorpdfstring{\pa{ਮਿਲਣਾ} · \pa{ਅਸੀਂ} \pa{ਮਿਲਦੇ} … (milnaa; asiin …)}{ਮਿਲਣਾ · ਅਸੀਂ ਮਿਲਦੇ … (milnaa; asiin …)}]{A bounded meeting-family payoff}

\label{lesson:PA-C188-payoff}



\begin{quote}
Recall the action name, the regular-group form, and the completed masculine-singular form in that order.
\end{quote}

\subsection*{Your turn: four skills}
\begin{itemize}
\raggedright
  \item \emph{Hear:} \emph{milnaa; asiin milde haan; main Aman nuun miliaa}; YOU SAY: \textbf{\pa{ਮਿਲਣਾ}; \pa{ਅਸੀਂ} \pa{ਮਿਲਦੇ} \pa{ਹਾਂ}; \pa{ਮੈਂ} \pa{ਅਮਨ} \pa{ਨੂੰ} \pa{ਮਿਲਿਆ}}
  \item \emph{Read it:} \textbf{\pa{ਮਿਲਣਾ} · \pa{ਅਸੀਂ} \pa{ਮਿਲਦੇ} \pa{ਹਾਂ} · \pa{ਮੈਂ} \pa{ਅਮਨ} \pa{ਨੂੰ} \pa{ਮਿਲਿਆ}}
  \item \emph{Write it:} cover the line and write \textbf{\pa{ਮਿਲਣਾ} · \pa{ਮਿਲਦੇ} · \pa{ਮਿਲਿਆ}}
\end{itemize}

Check the three jobs: name an action, describe a regular group meeting, and report one completed masculine-singular meeting.

\begin{tcolorbox}[breakable,colback=teal!4,colframe=teal!35!black,title={Before you move on}]
Produce the three-way contrast once more without looking.

This is the eleventh prerequisite tranche exposed by Punjabi A2 mock drafting. It does not unblock Punjabi A2 mock 1 by itself. It is not an A2 mock, not A2 readiness or attainment, and not external or human validation.
\end{tcolorbox}
